\documentclass[11pt]{article}
\usepackage[margin=1in]{geometry}
\usepackage{times}
\usepackage{enumitem}
\usepackage[hidelinks]{hyperref}
\usepackage{amsmath}
\setlength{\parskip}{4pt}
\setlength{\parindent}{0pt}

\title{\textbf{Research Proposal}\\[4pt]
\large Dynamic Partial Repartitioning of DAG Workflows After
Resource Failures in Federated Fog Environments}
\author{Gayani Gupta\\
Department of Computer Science and Engineering\\
University of North Texas}
\date{\today}

\begin{document}
\maketitle

\section{Introduction}

Workflows that respond to disasters, such as flood forecasting,
wildfire tracking, and evacuation planning, run as directed
acyclic graphs (DAGs) of dependent tasks. In a federated fog
system these tasks are placed on resources close to the event,
because the link to the cloud may be slow or unavailable exactly
when the workflow is needed most.

My previous work studied how to place these workflows well. Our
IEEE Cloud Summit 2025 paper defined the allocation problem and
compared four scheduling methods~\cite{gupta2025efficient}. My
current work extends that study: it makes the model data-aware,
adds a measurable accuracy bound, and tests the schedulers under
deadline slack, federation growth, workflow growth, task
failures, and link-bandwidth
degradation~\cite{gupta2026dataaware,gupta2026code}.

All of this work has one common assumption: the schedule is
decided once, before execution starts. In a real deployment this
is not enough. A fog resource can fail while the workflow is
running, and every unfinished task assigned to that resource is
stuck. Recomputing the complete schedule takes too much time,
moves tasks that are already running correctly, and increases
communication and migration costs.

This proposal presents a different method. When a resource
fails, we identify only the part of the task DAG that is affected
by the failure, repartition and reassign only that part to
available resources, and leave everything else unchanged.

\section{Background and my earlier work}

\subsection{Published work: IEEE Cloud Summit 2025}

In~\cite{gupta2025efficient} we introduced the task-allocation
problem for DAG workflows on federated fog systems. The model
connects task placement with data movement: if a task and its
predecessor run on different nodes, the data on the edge $D_{ji}$
must travel over a link with bandwidth $B_{ba}$, which costs
$D_{ji}/B_{ba}$ time. We built a mixed-integer linear program
(MILP) with correctness proofs, and three scalable methods: an LP
relaxation with randomized rounding, a data-aware greedy
scheduler, and a genetic algorithm. We compared them under
deadline slack, number of fog nodes, workflow size, and task
failure rate.

\subsection{Current work: data-aware scheduling}

In~\cite{gupta2026dataaware} I improved the same model and made
the experiments more complete. The work adds a maximum-accuracy
bound so each method can be checked against a clear reference; a
retry-once failure model, which separates accuracy from deadline
feasibility; a bandwidth-degradation experiment with a
bandwidth-blind baseline, which shows that a scheduler can look
good on accuracy while producing late schedules; and scalability
measurements that show where each method stops working (the MILP
near 800 assignment variables, the LP near $10^5$ constraints at
1000 tasks, and the genetic algorithm under a fixed search
budget). The main lesson is that the choice of scheduler matters
most when the deadline is tight, and the simple data-aware greedy
is the only method that is both fast and close to the bound at
1000 tasks.

The code, data, executed notebook, figures, and slides are public
at \url{https://github.com/gayanigupta/data-aware-scheduling}~\cite{gupta2026code}.

\subsection{What is still missing}

Both studies handle failures only at the task level: a failed
task is retried once on the same node. Neither study answers what
to do when a whole resource fails during execution. That is the
problem of this proposal.

\section{Research problem}

During workflow execution, the system monitors the fog resources.
When one resource fails, the unfinished tasks assigned to it can
no longer continue. A simple solution is to partition and
schedule the complete workflow again. However, this can take too
much time, it moves tasks that are already running correctly, and
it increases communication and migration costs.

I will develop a dynamic method that identifies the part of the
task DAG affected by the resource failure. The method
repartitions only that affected part and assigns its unfinished
tasks to suitable available fog resources. Tasks running on
unaffected resources remain unchanged.

\textbf{Main research question.} How can we quickly repartition
and reschedule only the affected part of a DAG workflow after a
fog resource fails, while reducing recovery time, communication
cost, task movement, and workflow completion time?

\section{Research steps}

\subsection{Represent the workflow}

I will represent the workflow as a task DAG. Each node is a task,
and each edge is a dependency and a data transfer between two
tasks.

\subsection{Represent the infrastructure}

I will represent the fog infrastructure as another graph. Each
node is a fog resource, and each edge is a network connection.

The resource information will include processing capacity,
available memory, current workload, and resource availability.
The network information will include communication delay,
available bandwidth, and data-transfer cost.

\subsection{Create the initial partition}

I will divide the task DAG into groups of connected tasks. These
groups are assigned to suitable fog resources based on execution
time, resource capacity, communication cost, and the workflow
deadline. For this step I can use the schedulers that already
exist in our simulator.

\subsection{Detect a resource failure}

During execution the system monitors the fog resources. When a
resource fails, the system identifies the failed node in the
infrastructure graph.

\subsection{Identify the affected tasks}

I will find the unfinished tasks assigned to the failed resource.
Completed tasks are not executed again. I will also check the
parent and child tasks of the affected tasks. This helps decide
whether moving only the failed tasks is enough, or whether some
dependent tasks must also be moved, for example to keep data
local or to still meet the deadline.

\subsection{Build the affected subgraph}

I will extract the affected tasks and their dependencies from the
complete task DAG. This smaller graph is the part that will be
repartitioned. Tasks running successfully on unaffected resources
stay where they are.

\subsection{Find replacement resources}

I will find fog resources that are still available and have
enough processing power and memory to run the affected tasks.
Shortest-path information will help identify resources that can
communicate well with the resources running the parent and child
tasks.

\subsection{Repartition the affected subgraph}

I will divide the affected subgraph into new task groups. The
decision considers task execution time, communication delay,
amount of transferred data, available resource capacity, the
workflow deadline, the cost of moving tasks, and load balance.

\subsection{Assign the new partitions}

I will assign the new task groups to the selected fog resources.
The goal is to recover from the failure while changing as little
of the original schedule as possible.

\subsection{Update communication paths}

After moving the tasks, I will update the communication paths
between the reassigned tasks and the tasks that stayed in place.
This part can build on the dynamic shortest-path idea of
Riazi, Bhowmick, and colleagues~\cite{riazi2018sssp,khanda2021sssp}.
Instead of computing all paths again, we update only the paths
affected by the failure and the reassignment.

\subsection{Continue the workflow}

The workflow continues from the last successful execution point.
Completed tasks are not restarted, and unaffected tasks are not
moved.

\section{Evaluation}

I will compare the proposed method against four baselines:
complete repartitioning and rescheduling; a greedy recovery
method that reassigns stranded tasks one by one; a static
schedule with no repair; and the MILP on small instances, as a
quality reference.

I will measure failure-recovery time, total workflow completion
time, communication cost, number of tasks moved, amount of data
transferred, energy use, deadline success rate, and the quality
of the repaired schedule.

I will also vary where the failure happens (a resource on the
critical path versus a peripheral resource), how many resources
fail at once, the workflow size, the federation size, and how
much slack is left at the time of the failure. Recovery is easy
when there is plenty of time left; the hard and interesting case
is when the deadline is nearly over.

\section{Novelty}

Several pieces of this problem already exist in the literature.
Riazi et al.~\cite{riazi2018sssp} and Khanda et
al.~\cite{khanda2021sssp} update a shortest-path tree after a
graph changes, but they do not apply this to repartitioning a
workflow after a resource failure. Yang et
al.~\cite{yang2024multiagent} study rescheduling in fog systems
when uncertain events occur. Lee et al.~\cite{lee2019fognetwork}
study how to redistribute tasks when fog resources join and
leave. Sakellariou and Zhao~\cite{sakellariou2004rescheduling}
reconsider parts of a workflow at carefully chosen points, and
classic schedulers such as HEFT~\cite{topcuoglu2002heft} assume
a static plan with full execution. None of these works connect a
failure in the infrastructure graph to the smallest affected
region of the task DAG.

The novelty of my work is the connection between the two graphs.
When a resource node fails, the method finds the unfinished tasks
on that resource, follows their dependencies in the task DAG,
identifies the smallest affected region, repartitions only that
region, finds replacement resources with enough capacity,
considers communication delay, data size, workload, and deadline,
updates only the affected communication paths, and continues the
workflow without restarting completed tasks.

To support this claim, the research must develop: a rule for the
boundary of the affected subgraph; a decision for whether
dependent tasks stay or move; the repartitioning algorithm; the
replacement-resource selection; the incremental path update; and
evidence that task movement and recovery time are reduced.

\section{Timeline}

\begin{itemize}[leftmargin=*]
\item \textbf{Months 1--3.} Extend the simulator to
resource-level failures. Implement the boundary rule and the
affected-subgraph extraction. Implement the complete-repartition
baseline.
\item \textbf{Months 4--6.} Build the repartitioning algorithm
and the replacement-resource selection. Add the incremental
path updates.
\item \textbf{Months 7--9.} Run the full comparison against the
baselines. Use the MILP as reference on small cases. Study how
sensitive the results are to the boundary rule.
\item \textbf{Months 10--12.} Write the paper, release the code
in the existing public repository~\cite{gupta2026code}, and
submit.
\end{itemize}

\section{Summary}

The three works form one line of research. The first paper built
the model and compared four ways to solve the static problem. The
current work made the model more complete and showed where each
method stops working. This proposal changes two things at once:
the failure granularity, from retrying a task to losing a whole
resource; and the scheduling granularity, from placing the whole
DAG to repairing only the smallest affected part. The platform,
the cost model, the schedulers, and the baselines already exist,
so the new work can build on them directly.

\begin{thebibliography}{9}\small
\bibitem{gupta2025efficient}
G.~Gupta, A.~Gupta, M.~Amini Salehi, A.~Fern\'andez Anta,
S.~Bhowmik, J.~Aguilar, and M.~Kourouma,
``Towards efficient allocation of tasks in DAG-based workflows
across federated fog systems,'' \emph{2025 IEEE Cloud Summit},
pp.~124--129, 2025.

\bibitem{gupta2026dataaware}
G.~Gupta, A.~Gupta, M.~Amini Salehi, A.~Fern\'andez Anta,
S.~Bhowmick, and J.~Aguilar,
``Data-aware scheduling of approximate DAG workflows in federated
fog systems,'' manuscript, 2026.

\bibitem{gupta2026code}
G.~Gupta, ``Data-aware scheduling of approximate DAG workflows in
federated fog systems: simulator and experiments,''
\url{https://github.com/gayanigupta/data-aware-scheduling}, 2026.

\bibitem{riazi2018sssp}
S.~Riazi, S.~Srinivasan, S.~K.~Das, S.~Bhowmick, and B.~Norris,
``Single-source shortest path tree for big dynamic graphs,''
\emph{2018 IEEE Int.\ Conf.\ on Big Data}, pp.~4054--4062, 2018.

\bibitem{khanda2021sssp}
A.~Khanda, S.~Srinivasan, S.~Bhowmick, B.~Norris, and S.~K.~Das,
``A parallel algorithm template for updating single-source
shortest paths in large-scale dynamic networks,'' \emph{IEEE
Trans.\ Parallel Distrib.\ Syst.}, vol.~33, no.~4, pp.~929--940,
2021.

\bibitem{yang2024multiagent}
Y.~Yang, F.~Ren, and M.~Zhang,
``A decentralized multiagent-based task scheduling framework for
handling uncertain events in fog computing,'' arXiv:2401.02219,
2024.

\bibitem{lee2019fognetwork}
G.~Lee, W.~Saad, and M.~Bennis,
``An online optimization framework for distributed fog network
formation with minimal latency,'' arXiv:1710.05239, 2019.

\bibitem{sakellariou2004rescheduling}
R.~Sakellariou and H.~Zhao,
``A low-cost rescheduling policy for efficient mapping of
workflows on grid systems,'' \emph{Scientific Programming},
vol.~12, no.~4, pp.~253--262, 2004.

\bibitem{topcuoglu2002heft}
H.~Topcuoglu, S.~Hariri, and M.-Y.~Wu,
``Performance-effective and low-complexity task scheduling for
heterogeneous computing,'' \emph{IEEE Trans.\ Parallel Distrib.\
Syst.}, vol.~13, no.~3, pp.~260--274, 2002.
\end{thebibliography}

\end{document}
